% ===================== 2. PROBLEM & MARKET =====================
\chapter{The problem and the market}

\section{How people currently travel}

Alexandria's shared transport is dominated by informal microbuses. They are cheap and frequent, and
they already establish the behaviour this service depends on: people accept sharing a vehicle,
accept walking to a boarding point, and accept a fixed fare for a known line.

What the informal system does not offer is a guaranteed seat, a known departure time, a recorded
payment, an identified driver, or any means of complaint. Those five absences are the whole of the
opportunity. This is a service-quality proposition, not a new mode of travel.

\section{Why universities first}

\begin{center}
\begin{tabularx}{\textwidth}{@{}L{40mm}Y@{}}
\toprule
\thead{Dense demand} & Shared transport requires demand density to function. Where demand is thin, vehicles run empty and the economics fail. \\
\thead{Predictable timing} & Students travel at the same hours on the same days, which is the condition scheduled routes are designed for. \\
\thead{Price sensitivity} & The population most likely to change behaviour for a lower fare. \\
\thead{Concentrated destinations} & Campus gates create natural route endpoints. \\
\thead{Word of mouth} & A closely connected population in which a service spreads or fails quickly and visibly. \\
\bottomrule
\end{tabularx}
\end{center}

\section{Market sizing}

\begin{notebox}
\notehead{Not claimed}
No total addressable market figure is presented. A credible bottom-up estimate requires corridor
demand data that has not been collected. Producing a top-down number from national population
statistics would not inform any decision, and is therefore omitted. The measurements required to
build a defensible estimate are listed in Chapter~7.
\end{notebox}

\section{Comparable operators}

\begin{center}
\begin{tabularx}{\textwidth}{@{}L{27mm}L{34mm}Y@{}}
\toprule
\thead{Operator} & \thead{Model} & \thead{Relevance} \\
\midrule
Swvl & Scheduled buses and vans, seat booking, fixed routes & Cairo-founded, closest comparable. Consumer mass-market was loss-making; profitability reported in 2025 on recurring B2B and B2G contracts. \\
Uber, Careem & On-demand private and pooled rides & Different price tier. Operate under Egypt's ride-hailing regulation. \\
Informal microbuses & Fixed lines, fixed fares, no booking & The actual incumbent and the real competitor. Cheap and frequent; no guarantee, no record, no accountability. \\
Via, Lyft (international) & Walk-to-corner shared rides & Source of the pricing evidence used in this design. Confirms walk-to-a-point sustains materially deeper discounts. \\
\bottomrule
\end{tabularx}
\end{center}

% ===================== 3. HOW IT WORKS =====================
\chapter{How the service works}

\section{The passenger}

\begin{enumerate}[label=\textbf{\arabic*.}]
\item \textbf{Choose a route.} Routes are named the way people already speak about them, for example
      Smouha to Alexandria University. The fare and frequency are shown alongside.
\item \textbf{Choose a boarding point.} Stops on the route are listed nearest first, with the
      walking time shown honestly. A more expensive option to be collected from the passenger's own
      street is offered alongside.
\item \textbf{Choose a departure.} At most three options are shown, each on a single line: departure
      time, boarding point, timetable arrival, price.
\item \textbf{Travel.} The passenger walks to the stop, boards by presenting a code, and leaves the
      vehicle at any point along the route.
\end{enumerate}

The interface never asks which product the passenger wants. It asks where they are going and when.
Internal vocabulary does not appear on screen.

\section{The driver}

A driver browses the published departure slots on approved routes and claims the ones they want.
Two steps. They do not design routes, set prices or forecast demand.

This division was chosen deliberately. Requiring drivers to plan routes and predict demand imposes
unpaid work before any earnings, which is a recognised cause of supply failure in peer-to-peer
transport. A study of an in-house microtransit service found that the flexibility drivers valued was
choosing \emph{shifts}, not designing service.

\section{The operator}

\begin{center}
\begin{tabularx}{\textwidth}{@{}L{44mm}Y@{}}
\toprule
\thead{Route network} & Which stops exist, in what order, at what fare, during which hours. \\
\thead{Departure grid} & Published slots, for example every fifteen minutes between 06:00 and 10:00. \\
\thead{Coverage} & Visibility of unclaimed slots where passengers are waiting, with the ability to offer a bonus for claiming one. \\
\thead{Verification} & Manual review of every driver: national identity document, driving licence, vehicle licence, photograph. \\
\thead{Support} & Complaint handling, incident response and an on-call safety line. \\
\bottomrule
\end{tabularx}
\end{center}

\section{Meeting points}

Every boarding point is created and verified by a person who physically visits the location. The
mandatory record is a name and a position; photographs, a night-safety flag and an accessibility
flag are optional and encouraged.

This matters for a specific reason. A 2017 study in \emph{PLOS ONE} measuring OpenStreetMap
completeness worldwide identified Egypt among the least complete, estimating that fewer than one
third of streets were mapped. The same study found that dense cities and main roads are the
best-mapped parts of any country. Because boarding points are surveyed by hand and routes follow
main corridors, the design depends very little on the weakest map data. Verifying this for the
launch corridors is an engineering task listed in Chapter~7.

% ===================== 4. REVENUE =====================
\chapter{Revenue and cost mechanics}

\begin{notebox}
\notehead{No projections}
This chapter describes how money is intended to move. It contains no forecast of volume, revenue or
profit. The service has never operated, so any such figure would be invention.
\end{notebox}

\section{Where revenue comes from}

\begin{center}
\begin{tabularx}{\textwidth}{@{}L{34mm}Y@{}}
\toprule
\thead{Commission} & A percentage of each fare. The default model. Rate is a configuration value set per city, route or driver segment. \\
\thead{Subscriptions} & Weekly or monthly commuter passes with a guaranteed seat. Predictable revenue and the strongest retention mechanism. \\
\thead{Upgraded boarding} & A higher-priced ticket for collection from the passenger's own street rather than a stop. \\
\thead{Alternative models} & Driver subscription, or a hybrid, are supported as configuration and can be enabled without engineering work. \\
\bottomrule
\end{tabularx}
\end{center}

\section{How a fare is set}

The stop-boarding fare is one flat figure per route. The street-collection fare is calculated in
three layers: a fixed uplift, a charge banded by the size of the detour, and a cost-recovery
component applied only if the first two do not cover the measured cost. The result is clamped
between a floor and a ceiling, and it is shown to the passenger before they commit.

Beyond a defined detour threshold, street collection is not expensive --- it is unavailable. A
deviation that would delay the passengers already aboard is refused rather than priced.

\section{Payment}

Wallet balance is the primary method. Card, mobile wallet, kiosk cash-in, bank transfer and Apple
Pay are supported for topping up or for direct payment, through a single payment provider
integration.

Cash is accepted, with one rule: it is always recorded against the scanned booking. The driver marks
the fare as collected, the amount becomes a liability owed by that driver to the platform, and it is
deducted from their next payout. There is no untracked cash sale. A driver whose outstanding cash
liability exceeds a configured limit is automatically prevented from accepting further cash
bookings, having been warned in advance.

\section{Cost structure}

\begin{center}
\begin{tabularx}{\textwidth}{@{}L{40mm}Y@{}}
\toprule
\thead{Routing and mapping} & Self-hosted routing on open map data. Published comparisons put a self-hosted engine at roughly \$20 per month against roughly \$510 per day for equivalent commercial matrix queries at scale. One team reported moving from about \$8{,}000 per month to about \$520. \\
\thead{Infrastructure} & Self-managed servers rather than managed cloud services, with a database replica and off-site backups. \\
\thead{Payment processing} & Per-transaction and per-payout fees. \\
\thead{Messaging} & Verification codes by SMS. All other messages are in-application, which avoids per-message cost. \\
\thead{Field operations} & Surveying and verifying boarding points, and reviewing driver documents. Labour, not software. \\
\bottomrule
\end{tabularx}
\end{center}

\begin{notebox}
\notehead{Not yet quoted}
Infrastructure figures above are drawn from published third-party comparisons, not from quotations
obtained for this project. Actual costs have not been established.
\end{notebox}
